%! Unit 2: Solving Linear Equations Ax = b — Summative Assessment — Unit Test (blank)
\documentclass[10pt]{article}
\usepackage{linalg-article}
\usepackage{linalg-boxes}
\newcommand{\bb}{\mathbf{b}}
\newcommand{\xx}{\mathbf{x}}
\newcommand{\cc}{\mathbf{c}}

% Part-divider strip
\newcommand{\parthead}[1]{%
  \vspace{6pt}\noindent
  \begin{headlinebox}{burgundy}{\color{white}\bfseries #1}\end{headlinebox}%
  \vspace{4pt}}

\begin{document}

\pageheader{Unit 2: Solving Linear Equations Ax = b}{Unit Test}
\namedateperiod

\vspace{0.06in}
\noindent\textbf{Instructions:} Show all work. No notes unless your teacher says so.

% ======================================================================
\parthead{Part A --- Vocabulary (8 pts)}
Write the letter of the best description next to each term.

\vspace{4pt}
\noindent
\begin{minipage}[t]{0.44\textwidth}
\begin{enumerate}[leftmargin=2.2em, itemsep=7pt, label=\arabic*.]
  \item $LU$ factorization \blank{1.2cm}
  \item Transpose \blank{1.2cm}
  \item Multiplier \blank{1.2cm}
  \item Singular matrix \blank{1.2cm}
  \item Pivot \blank{1.2cm}
  \item Permutation matrix \blank{1.2cm}
  \item Inverse matrix \blank{1.2cm}
  \item Upper triangular ($U$) \blank{1.2cm}
\end{enumerate}
\end{minipage}%
\hfill
\begin{minipage}[t]{0.53\textwidth}
\small
\begin{description}[leftmargin=1.4em, itemsep=3pt]
  \item[(A)] $A^{-1}$, the matrix that undoes $A$: $A^{-1}A=I$.
  \item[(B)] The matrix $A^{\mathsf T}$ formed by flipping $A$ across its diagonal (rows $\leftrightarrow$ columns).
  \item[(C)] The identity with two rows swapped; it records a row exchange.
  \item[(D)] The first nonzero entry in a row, used to clear the entries below it.
  \item[(E)] A matrix with all zeros below the main diagonal.
  \item[(F)] $A=LU$: a lower-triangular $L$ times an upper-triangular $U$.
  \item[(G)] A matrix with no inverse --- its determinant $ad-bc$ is $0$.
  \item[(H)] $\ell=(\text{entry to clear})\div(\text{pivot})$; how many pivot rows you subtract.
\end{description}
\end{minipage}

% ======================================================================
\parthead{Part B --- Multiple Choice (12 pts)}
Circle the best answer.

\begin{enumerate}[leftmargin=1.9em, itemsep=9pt]
  \item In $A=LU$, the pivots of the elimination appear on the diagonal of
    \hfill (A) $L$ \quad (B) $U$ \quad (C) $A^{-1}$ \quad (D) the permutation $P$
  \item To clear the entry below a pivot $p$, the multiplier you subtract is $\ell=$
    \hfill (A) $(\text{entry})\div p$ \quad (B) $p\times(\text{entry})$ \quad (C) $p+(\text{entry})$ \quad (D) $p\div(\text{entry})$
  \item If $P$ is a permutation matrix, then $P^{\mathsf T}$ equals
    \hfill (A) $P^2$ \quad (B) $-P$ \quad (C) $P^{-1}$ \quad (D) the identity $I$
  \item A $2\times2$ system $A\xx=\bb$ has \textbf{exactly one} solution when the two lines
    \hfill (A) are the same line \quad (B) cross at one point \quad (C) are parallel but different \quad (D) are both horizontal
  \item For any two matrices that can be multiplied, $(AB)^{\mathsf T}=$
    \hfill (A) $A^{\mathsf T}B^{\mathsf T}$ \quad (B) $B^{\mathsf T}A^{\mathsf T}$ \quad (C) $AB$ \quad (D) $BA$
  \item A square matrix $A$ has \textbf{no} inverse exactly when elimination produces
    \hfill (A) a zero pivot \quad (B) all $1$'s \quad (C) a symmetric $U$ \quad (D) a row swap
\end{enumerate}

% ======================================================================
\parthead{Part C --- Short Answer \& Computation (30 pts)}
Show your work.

\begin{enumerate}[leftmargin=1.9em, itemsep=13pt]
  \item Solve the system by \textbf{elimination}, then back-substitute and check. State the pivot and the
        multiplier you used.
        \[ 2x+y=7, \qquad 4x+3y=17. \]
        \vspace{1.8cm}
  \item Row-reduce $A=\begin{bsmallmatrix}1&3&1\\3&11&5\\2&14&11\end{bsmallmatrix}$ to upper-triangular $U$.
        List the three multipliers and the pivots.
        \vspace{2.4cm}
  \item Let $A=\begin{bsmallmatrix}3&1\\5&2\end{bsmallmatrix}$. Use the $ad-bc$ formula to find $A^{-1}$,
        then check that $A^{-1}A=I$.
        \vspace{1.8cm}
  \item Using the inverse from the previous problem, solve $A\xx=\bb$ for
        $\bb=\begin{bsmallmatrix}4\\7\end{bsmallmatrix}$ (compute $\xx=A^{-1}\bb$).
        \vspace{1.6cm}
  \item Factor $A=\begin{bsmallmatrix}3&6\\9&20\end{bsmallmatrix}$ as $A=LU$. Write $L$ and $U$, then verify
        $LU=A$.
        \vspace{1.8cm}
  \item Using the $L$ and $U$ from the previous problem, solve $A\xx=\bb$ for
        $\bb=\begin{bsmallmatrix}3\\11\end{bsmallmatrix}$ in \textbf{two sweeps}: first $L\cc=\bb$
        (forward), then $U\xx=\cc$ (back).
        \vspace{1.8cm}
  \item \textbf{(a)} The matrix $A=\begin{bsmallmatrix}0&3\\2&5\end{bsmallmatrix}$ has a $0$ in the top-left
        pivot spot. Write the permutation matrix $P$ that swaps the two rows, and give $PA$.
        \textbf{(b)} Write the transpose of $B=\begin{bsmallmatrix}4&1\\1&3\end{bsmallmatrix}$ and state
        whether $B$ is symmetric.
        \vspace{1.8cm}
\end{enumerate}

% ======================================================================
\parthead{Part D --- Extended Response (10 pts)}
Explain your reasoning in full sentences.

\begin{enumerate}[leftmargin=1.9em, itemsep=16pt]
  \item Let $A=\begin{bsmallmatrix}1&3\\2&6\end{bsmallmatrix}$.
        \textbf{(a)} Compute $ad-bc$. Does $A$ have an inverse? What does that say about solving $A\xx=\bb$?
        \textbf{(b)} How many solutions does $A\xx=\begin{bsmallmatrix}4\\8\end{bsmallmatrix}$ have?
        \textbf{(c)} How many does $A\xx=\begin{bsmallmatrix}4\\7\end{bsmallmatrix}$ have? Justify (b) and (c)
        using elimination or the row/line picture.
        \vspace{2.6cm}
  \item Let $A=\begin{bsmallmatrix}1&4\\2&1\end{bsmallmatrix}$.
        \textbf{(a)} Compute $A^{\mathsf T}A$ and check that it is symmetric.
        \textbf{(b)} Explain why $A^{\mathsf T}A$ is \emph{always} symmetric for any matrix $A$, using the
        reversal rule $(AB)^{\mathsf T}=B^{\mathsf T}A^{\mathsf T}$.
        \vspace{2.4cm}
\end{enumerate}

\end{document}
